\documentclass[12pt, a4paper]{article}

\usepackage[margin=1.5in]{geometry}
\usepackage{graphicx}
\usepackage{amsmath}
\usepackage{amssymb}
\usepackage{booktabs}
\usepackage{hyperref}
\usepackage{threeparttable}
\usepackage{caption}
\usepackage{subcaption}
\usepackage[utf8]{inputenc}
\usepackage[T1]{fontenc}
\usepackage{natbib}
\usepackage{doi}
\usepackage{CJKutf8}

\title{Forecast Dispersion and Forecast Errors across Firms and Time}
\author{Hatsu Kawabata \and Tatsuro Senga}
\date{May 29, 2024}


\begin{document}
\begin{CJK}{UTF8}{min}
\maketitle

\begin{abstract}
This paper investigates the properties of analyst forecast dispersion and forecast errors using a comprehensive dataset of Japanese firms from 1985 to 2023. We construct time-series indices of forecast dispersion and errors and explore their relationships with macroeconomic and financial market indicators. Our analysis reveals that forecast dispersion and errors are positively correlated, indicating that greater disagreement among analysts is associated with larger forecast errors. Forecast dispersion tends to be smaller for larger firms with more analyst coverage, and the number of analysts covering a firm is positively related to its size, age, and stock volatility. We find that the forecast dispersion and error indices are correlated with other popular uncertainty measures like the Economic Policy Uncertainty index, with spikes corresponding to events that heightened uncertainty. The indices are also countercyclical and negatively correlated with stock market performance. Our findings highlight the role of firm-level uncertainty in macroeconomic fluctuations and demonstrate the usefulness of analyst forecast data in studying the relationship between information, uncertainty, and the macroeconomy.
\end{abstract}

\section{Introduction}
Expectations play a central role in macroeconomics and finance, influencing the decisions of households, businesses, and policymakers.
Yet, empirical studies testing the theoretical implications and establishing the stylized facts on expectations and more fundamentally
on how economic agents acquire and use information are still limited, primarily because of the limited availability of data to use.
In this paper, we study a relatively underused dataset on market analysts' forecasts of earnings of firms called IBES (Institutional Brokers' Estimate System)
to show both cross-sectional and time-series patterns in analyst forecasts.
Focusing on Japanese publicly traded firms, we use data taken from IBES and merge it with NikkeiNeeds and DBJ data to construct time-series indices
of forecast dispersion and errors, which provide a real-time measure of micro-level uncertainty.
We then explore the relationships between these indices and various firm characteristics, as well as macroeconomic and financial market indicators.

Our analysis reveals several key findings.
First, forecast dispersion and errors are positively correlated, indicating that greater disagreement among analysts
is associated with larger forecast errors.
This suggests that dispersion reflects genuine uncertainty rather than differences in opinion or information.

Second, the number of analysts covering a firm plays an important role in shaping forecast properties.
Dispersion and errors tend to be smaller for firms with more analyst coverage,
even after controlling for other factors such as firm size, age, and volatility.
We interpret this as evidence of information spillovers and learning among analysts. Having more analysts covering a firm allows for greater information sharing and convergence of views, leading to reduced dispersion and more accurate consensus forecasts. This highlights the importance of analyst coverage in reducing firm-level uncertainty.

Third, the forecast dispersion and error indices exhibit significant time-series variation,
with spikes corresponding to major economic and geopolitical events that heighten uncertainty.
The indices are positively correlated with other popular uncertainty measures like
the Economic Policy Uncertainty (EPU) index for Japan. They also display a clear countercyclical pattern,
rising during downturns and periods of stock market weakness. This countercyclical behavior suggests that firm-level uncertainty rises during bad economic times, potentially amplifying macroeconomic fluctuations.

Overall, our findings highlight the usefulness of analyst forecast data
in measuring and understanding the evolution of firm-level uncertainty over time.
The indices we develop can serve as a valuable tool for policymakers and market participants
to monitor uncertainty in real time. By capturing the degree of disagreement and inaccuracy in analyst expectations, these indices provide a novel, forward-looking measure of the uncertainty surrounding firms' fundamentals and prospects. As such, they complement existing macro uncertainty measures and can help improve our understanding of the links between micro-level information frictions and aggregate outcomes.

The dispersion of analyst forecasts has been a topic of interest in the accounting literature, while 
economic analysis is yet still scarce.\footnote{See \citet{NaraNoma2024} for a recent survey and related papers referenced therein.}
The dispersion of analyst forecasts has been a topic of interest in various studies, which have examined its relationship with analyst coverage, accuracy and optimism of analyst forecasts, corporate disclosure, and firm characteristics. \citet{Diether2002}, \citet{Ciccone2003}, and \citet{LiuNatarajan2012} have used the dispersion of analyst forecasts as a proxy for future uncertainty and found that it is positively correlated with analyst coverage, indicating that higher future uncertainty leads to greater analyst coverage.
Regarding firm size, \citet{LangLundholm1996}, \citet{Diether2002}, and \citet{LiuNatarajan2012} have shown that larger firms tend to have lower dispersion of analyst forecasts. \citet{Clement2003} report that the disclosure of management forecasts leads to a reduction in the dispersion of analyst forecasts, suggesting that management forecasts have the effect of reducing future uncertainty.
\citet{LangLundholm1996} also find that firms with higher disclosure ratings have lower dispersion of analyst forecasts. Their study reveals that firms with higher standard deviation of ROE, higher correlation between stock prices and earnings, larger earnings surprises, and more recent forecast revisions tend to have higher dispersion of analyst forecasts.
Furthermore, \citet{Brown1987a} state that there is a positive relationship between the accuracy of analyst forecasts and the dispersion of analyst forecasts.
In the Japanese context, while there is no study that directly analyzes the relationship between the dispersion of analyst forecasts and the value relevance of analyst forecasts, \citet{Obinata2010} has examined the relationship between the dispersion of analyst forecasts and the value relevance of actual earnings.

Our paper contributes to several strands of literature in Economics.
Firstly, we add to the research on the properties and determinants of analyst forecasts. 
Prior studies have documented cross-sectional dispersion in analyst forecasts and its relationship with other firm characteristics \cite{hong2000bad, brennan2000investment, diether2002differences}, 
as well as the time-series dynamics of aggregate forecast dispersion \cite{yu2011analyst, barinov2022firm}. 
We extend this literature by conducting a systematic analysis of both forecast dispersion and errors in the context of Japanese firms, 
examining their relationships with an extensive set of firm and macro variables. 
Our results corroborate some of the key findings in the US while also providing new insights.

Secondly, we contribute to the empirical literature on uncertainty. Several proxies have been de-
veloped within the literature, ranging from the volatility of GDP or stock prices to disagreement
and forecast errors in survey data, as uncertainty is difficult to identify. For example, \cite{leahy1996investment} construct a measure of uncertainty from the volatility of stock returns for indi-
vidual firms. \cite{guiso1999investment} use survey data on demand forecasts by Italian firms to
infer the level of uncertainty facing these individual firms.\footnote{Guiso and Parigi (1999) use 3-point probability distributions from the Bank of Italy Survey of Investment
in manufacturing (SIM) and Morikawa (2013) uses 2-point distributions from his original survey and finds that
uncertainty related to the tax system and trade policy matters for firms' capital investment and overseas activities.}
\cite{bond2005uncertainty} consider several uncertainty measures based on analyst forecast data.
In this paper, we use data on earnings forecasts by individual analysts as in
\cite{bond2005uncertainty}; however, we examine not only the average cross-sectional distribution but
also the cyclical changes of this uncertainty measure. \cite{bachmann2013uncertainty} use survey data
from the IFO Business Climate Survey, which asks forecasters about their own future prospects
rather than about macroeconomic variables such as GDP, to extensively study various measures
of uncertainty. We also use forecast disagreement to measure uncertainty, building on these studies.

Furthermore, our indices capture uncertainty at the firm level, complementing the growing literature on macroeconomic and policy uncertainty \citep{bloom2009impact, baker2016measuring, jurado2015measuring}. By aggregating firm-level forecast dispersion and errors into time-series indices, we provide novel measures of the average uncertainty faced by firms over time. This allows us to study the relationship between micro-level uncertainty and macro aggregates, shedding light on the transmission and amplification of uncertainty shocks.

The rest of the paper is organized as follows. Section 2 describes the data and methodology. Section 3 presents our empirical findings on the cross-sectional determinants of forecast dispersion and errors. Section 4 introduces the time-series indices and explores their properties and correlations. Finally, Section 5 concludes.



\bigskip
\section{Data and Methodology}
We construct a comprehensive dataset by combining earnings forecasts made by analysts from IBES, stock prices from NikkeiNeeds, and financial information from the Development Bank of Japan (DBJ). 
We define key variables such as forecast dispersion, forecast errors, and the number of analysts covering each firm.
Our dataset covers a wide range of public firms in Japan from 1985 to 2023, allowing us to study the evolution of forecast dispersion and errors over time. 

\subsection{IBES Data}
IBES Data is a database that collects and compiles financial analyst earnings estimates and recommendations. 
It provides a comprehensive view of analyst expectations and consensus estimates, including earnings per share (EPS) estimates, buy/hold/sell recommendations, and price targets. 
The data covers a wide range of companies across various industries and regions and is available at different frequencies.

Figure \ref{fig:merge_nikkei_counts} shows the number of firms in IBES matched with NikkeiNeeds and DBJ data over time.
The fraction of firms in NikkeiNeeds and DBJ data being matched with IBES data was just below 30\% at early stage (450 firms out of 1,500 firms),
while the fraction started to increase to 50\% at the end of 2010 (1,000 firms out of 2,000 firms).
As become clear later, the matched sample and the unmatched sample have similar characteristics in terms of the dispersion of forecasts and forecast errors.



\begin{figure}[!ht]
    \centering
    \includegraphics[width=1.0\textwidth]{graph/merge_nikkei_counts.png}
    \caption{IBES being matched with NikkeiNeeds and DBJ}
    \label{fig:merge_nikkei_counts}
\end{figure}



\subsection{Dataset Construction}
For the analysis below, we construct two datasets: Dataset A combines IBES and NikkeiNeeds data, including consensus forecasts, forecast dispersion, forecast errors, stock prices and volatility at a monthly frequency from January 1985 to June 2023.
Dataset B adds financial information from the Development Bank of Japan to Dataset A at an annual frequency from 1985 to 2023.

We define key variables such as forecast dispersion, forecast errors, and the number of analysts covering each firm. 
To this end, we take the earnings per share (EPS) forecasts from IBES, $f_{hijt}$, where $h$ represents the forecast horizon, such as six months ahead of the fiscal year end. 
The index $i$ denotes the analyst making the forecast. 
The index $j$ refers to the firm being forecasted, and $t$ indicates the forecast time, for example, January 2005.

Forecast dispersion is calculated as:
\begin{equation}
    Fdis_{hjt} = \frac{\sigma_{hjt}}{|\bar{f}_{hjt}|} = \frac{\sqrt{\frac{1}{N_{hjt}-1} \sum_{i=1}^{N_{hjt}} (f_{hijt} - \bar{f}_{hjt})^2}}{|\bar{f}_{hjt}|}
\end{equation}
where $N_{hjt}$ is the number of analysts making forecasts for the given forecast horizon $h$, firm $j$, and time $t$, and $\bar{f}_{hjt}$ is the median of the forecasts across analysts for the given $h$, $j$, and $t$.

Forecast errors are calculated as:
\begin{equation}
    FE_{hjt} = \left| \log \left( \frac{e_{jy(t)}}{\bar{f}_{hjt}} \right) \right|
\end{equation}
where $e_{jy(t)}$ is the realized earnings per share for firm $j$ and year $y(t)$ of time $t$. For instance, $y(t)$ for $t=(\text{2004 January})$ is 2004.



\subsection{Descriptive Statistics}

\begin{table}[htbp]
\centering
\resizebox{1.\textwidth}{!}{
\begin{threeparttable}
\input{table/desc_stats_a.tex}
\begin{tablenotes}
\footnotesize
\item
\end{tablenotes}
\end{threeparttable}
}
\caption{Descriptive Statistics: Dataset A}
\label{tab:desc_stats_a}
\end{table}

Table \ref{tab:desc_stats_a} presents descriptive statistics for key variables in Dataset A.
The mean realized EPS is 71.4 Japanese yen, while the mean median estimated EPS is 89.8 Japanese yen.
While the median estimated EPS tends to be higher than realized EPS across the percentile distribution,
the 5th percentile estimated EPS (minus 9.8 Japanese yen) is exceptionally much higher than the 5th percentile realized EPS (minus 94.2 Japanese yen).
This is related to the fact that the standard deviation of the estimated EPS is smaller than the standard deviation of the realized EPS.
The fact that the standard deviation of the estimated EPS is smaller than the standard deviation of the realized EPS is consistent with 
a simple model wherein analysts form expectations about future earnings of each firm under imperfect information.
A simple Bayes rule implies that, suppose the prior mean of analysts is around the unconditional mean or median, then the analyst forecast distribution 
tends to be more compressed around the prior mean because an analyst's forecast is the weighted average of the prior mean and a signal 
which is around a true value of earnings of each firm in the future. 
In fact, a similar fact about the standard deviation of forecasts has been found in other contexts like sales forecasts of business managers. 
The distribution of business managers about their own sales growth tend to be more compressed than that of 
realized sales growth.\footnote{See Bloom st al. (2021) and Barrero (2022) for more details.}

Table 1 also shows how many analysts cover each firm. 
The mean number of analysts covering each firm is 4.8 and the median is 3 in the sample. 
The 95th percentile of the number of estimates is 15, and the smallest by definition in this IBES data includes firms that have only one analyst covering.
For the analysis below, we focus on firms that have at least two analysts covering each firm and the results are robust to this restriction as become clear 
later.\footnote{The main results remain unchanged when we focus on the sample of firms with at least five analysts or ten analysts covering each firm.}

\begin{table}[htbp]
\centering
\resizebox{1.0\textwidth}{!}{
\begin{threeparttable}
\input{table/desc_stats_b.tex}
\begin{tablenotes}
\footnotesize
\item
\end{tablenotes}
\end{threeparttable}
}
\caption{Descriptive Statistics: Dataset B}
\label{tab:desc_stats_b}
\end{table}


Table 2 shows descriptive statistics for key variables in Dataset B, which includes financial information taken from the DBJ data. 
First, we can compare realized EPS, median EPS, number of estimates across two datasets A and B. 
It's clear that whether or not we can match with the DBJ data doesn't materially change the characteristics of the underlying datasets. 
This is also true for forecast dispersion and forecast error measures created across the two datasets. 
Since financial information is included in Dataset B, we can obtain the summary statistics in terms of market capitalization and sales for Dataset B. 
For instance, the mean sales is 420 billion Japanese yen and the mean market capitalization is 270 billion Japanese yen.




Table \ref{tab:desc_stats_b} shows descriptive statistics for Dataset B, which includes financial variables.
Table 2 shows descriptive statistics for key variables in Dataset B, which includes financial information taken from the DBJ data. 
First, we can compare realized EPS, median EPS, number of estimates across those two datasets. 
It's clear that whether or not we can match with the DBJ data doesn't materially change the characteristics of the underlying datasets. 
This is also true for forecast dispersion and forecast error measures created across the two datasets. 
Since financial information is included in Dataset B, we can gauge the summary statistics in terms of market capitalization and sales. 
For instance, the mean sales is 420 billion Japanese yen and the mean market capitalization is 270 billion Japanese yen.


\bigskip
\section{Empirical Findings}
\subsection{Timing and Horizon of Forecasts}
The constructed data includes analysts' forecasts of EPS for each firm's accounting year, with each forecast being made at monthly frequency, ranging from Horizon 0 to Horizon 12 and even further.
Horizon 0 corresponds to the fiscal year end and analysts still can report their forecasts for the fiscal year end and they do indeed.
In this subsection, we therefore look at how forecasts are released and the relationship between the timing of the forecasts and the accuracy and dispersion of the forecasts.
Figure \ref{fig:numfirm_h} shows that the majority of analyst forecasts start to be released 10 months ahead of the fiscal year end.
By the time we reach 9 months ahead of the fiscal year, there are almost 100\% of forecasts available for all the firms in the dataset, 
and after that, there is a stable number of firms that have forecasts available until the end of the fiscal year.


\begin{figure}[htbp]
\centering
\includegraphics[width=1.0\textwidth]{graph/numfirm_h.png}
\caption{Number of Firms with Forecasts by Forecast Horizon}
\label{fig:numfirm_h}
\end{figure}

Focusing on the relationship between the timing of the forecasts and the accuracy and dispersion of the forecasts for Horizon 0 to Horizon 8 or 9 when most of the forecasts are released, 
Figure \ref{fig:fdis_fe_h} shows the mean forecast dispersion and forecast errors for each horizon tend to decrease as the forecast horizon approaches the fiscal year end.
Looking at the forecast dispersion on the left panel of Figure \ref{fig:fdis_fe_h}, 
the degree of dispersion and hence disagreement among analysts peaks around Horizon 9 and then starts to decline. 
But the speed of such convergence of forecasts among analysts seems to be slower than the pace of getting better informed, 
as seen in the right panel of Figure \ref{fig:fdis_fe_h}, where we can see that the size of forecast error decreases substantially from Horizon 8 towards the end of the fiscal year, which is Horizon 0.

\begin{figure}[htbp]
\centering
\includegraphics[width=0.47\textwidth]{graph/meanfdis_h.png}
\includegraphics[width=0.47\textwidth]{graph/meanFElog_h.png}
\caption{Forecast Dispersion and Errors over Forecast Horizon}
\label{fig:fdis_fe_h}
\end{figure}



\subsection{Number of Forecasts}

In the previous section, we have seen that the dispersion of analysts' forecasts decreases as the forecast horizon approaches the fiscal year end. 
We also have seen that the accuracy of analysts' forecasts increases as the forecast horizon approaches the fiscal year end.
In stead of looking at the within-firm variation of the dispersion and accuracy of forecasts, this section explores the between-firm variation of the dispersion and accuracy of forecasts, 
by looking at various firm characteristics including the number of analysts, firm size, and firm age.

The left panel of Figure \ref{fig:sd_actual_numest} shows the relationship between the number of estimators and the standard deviation of forecasts. 
As the panel shows, the standard deviation of EPS forecasts is larger when the number of forecasters is large. 
To explore more this, the right panel of Figure \ref{fig:sd_actual_numest} shows the relationship between the realized level of EPS for each firm and the number of forecasters covering that firm. 
As seen in the panel, that relationship is positive - the larger the realized level of earnings (EPS), the more the number of coverage analysts. 
The positive relationship between the standard deviation of analyst forecasts of EPS and the number of estimators in the left panel might simply be a mechanical result of the positive correlation between the level of realized EPS and the number of estimators, as seen in the right panel.

\begin{figure}[htbp]
\centering
\includegraphics[width=0.45\textwidth]{graph/StdevNumest.png}
\includegraphics[width=0.45\textwidth]{graph/ActualNmest.png}
\caption{Standard Deviation of Forecasts and Realized EPS by Number of Estimates}
\label{fig:sd_actual_numest}
\end{figure}

To explore the relationship between the number of analysts and the accuracy and dispersion of the forecasts, 
Figure \ref{fig:fdis_fe_numest} shows the relationship between the coefficient of variation of the forecasts by analysts and the number of analysts. 
This way we control for the fact that the standard deviation might be larger just because the level of realized EPS is larger. 
As can be seen from the panel, the relationship between the coefficient of variation of the forecasts and the number of estimators is negative. 
To further explore this, the right panel of Figure \ref{fig:fdis_fe_numest} plots the relationship between forecast error and the number of estimates. 
The panel shows a negative relationship between the size of the forecast error made by the consensus forecast and the number of estimators. 

Combining the left and right panels, one can think of the mean forecast EPS is the result of aggregating each individual analyst's forecasts. 
As long as each analyst is partially informed about the firm's future earnings, the more analysts that cover each firm, the more informed the average forecast is and the less dispersed the forecasts are.
It is not necessary for this story to be true though, there may be information spillovers across analysts.
Therefore, the more analysts that cover each firm, we might expect that they learn from each other, information spills over, and the accuracy of information held by the analysts increases, leading to a lower degree of disagreement. 
To dig deeper into this point, we investigate the drivers of forecast dispersion and forecast error utilizing the panel feature of the dataset and the available control variables and fixed effects.

\begin{figure}[htbp]
\centering
\includegraphics[width=0.45\textwidth]{graph/FdisCVnumest.png}
\includegraphics[width=0.45\textwidth]{graph/FElogNumest.png}
\caption{Forecast Dispersion and Errors by Number of Estimates}
\label{fig:fdis_fe_numest}
\end{figure}




\subsection{Forecast dispersion is smaller for firms with more analysts}

The previous section considered the bivariate correlations; though, this section take this further by adding more variables and controlling for other factors in Table \ref{tab:reg_fdis}.
It presents regression results for the determinants of forecast dispersion, including the number of analysts, firm size, firm age, earnings volatility, and stock price volatility.
Column (1) reports the result of regressing a measure of forecast dispersion, confirming the statistical significance of the relationship in the right panel of Figure \ref{fig:fdis_fe_numest}.
Column (1) shows that a one unit increase in the number of estimates (NUMEST) is associated with a 0.014 decrease in forecast dispersion, suggesting that having more analysts reduces forecast dispersion.
Column (2) adds year and firm dummies, and column (3) adds the standard other controls such as firm size and firm age.
For firms with more analysts covering and reporting earnings forecasts, those analysts still make smaller forecast dispersion.

\begin{table}[htbp]

\centering
\resizebox{1.0\textwidth}{!}{
\begin{threeparttable}
\input{table/reg_T01.tex}
\begin{tablenotes}
\footnotesize
\end{tablenotes}
\end{threeparttable}
}
\caption{Regression Results for Forecast Dispersion}
\label{tab:reg_fdis}
\end{table}

Column (4) considers earnings volatility and stock price volatility as drivers of forecast dispersion.
These variables are natural candidates for drivers of forecast dispersion because it takes time to get better informed if one learns about volatile outcomes.
It may be less direct, however, the volatile stock returns may imply that the volatility of the earnings is higher.
As predicted, both earnings and stock price volatility are positively correlated with forecast dispersion, indicating analysts disagree more about firms that have more volatile earnings growth and stock prices.
Column (4) indicates that a one unit increase in the standard deviation of realized EPS is associated with a 0.035 increase in forecast dispersion.
It also shows that a one unit increase in the standard deviation of stock returns is associated with a 0.132 increase in forecast dispersion.
The main result in terms of the number of analysts remains robust and its magnitude is stonger in that adding more analysts to a firm reduces forecast dispersion by 0.118, 
while the mean value of the dependent variable is 0.2 as in Table \ref{tab:desc_stats_b}.

So far, we include all firms that have at least two analysts covering and reporting earnings forecasts.
In column (5), we exclude firms that have only two analysts covering and reporting earnings forecasts, and the result is robust to doing so.
Column (5) shows that a one unit increase in the number of analysts is associated with a 0.131 decrease in forecast dispersion, a similar magnitude as seen in column (4) for all firms.
Column (6) excludes firms that have fewer than 6 analysts covering and reporting earnings forecasts.
The results remain robust with such sample selection, and it is confirmed that the number of analysts drives the size of forecast dispersion.
In summary, the regression table examines factors associated with forecast dispersion, finding that having more estimates reduces dispersion, while greater variability in actual values and higher stock volatility are linked to more dispersed forecasts. The results are generally consistent across sub-sample specifications.





\subsection{Forecast errors are smaller for firms with more analysts}

Table \ref{tab:reg_fe} shows regression results for the determinants of forecast errors, including the number of analysts, firm size, firm age, earnings volatility, and stock price volatility.
Column (1) reports the result of regressing a measure of forecast error, confirming the statistical significance of the relationship in the right panel of Figure \ref{fig:fdis_fe_numest}.
Column (1) shows that a one unit increase in the number of estimates (NUMEST) is associated with a 0.011 decrease in forecast error, suggesting that having more estimates reduces forecast error.
Column (2) adds year and firm dummies, and column (3) adds the standard other controls such as firm size and firm age.
For firms with more analysts covering and reporting earnings forecasts, those analysts still make smaller forecast errors.
With more controls, column (3) indicates that adding an analyst to a firm is associated with a 0.141 decrease in forecast error, a larger magnitude compared to column (1).

\begin{table}[htbp]
\centering
\resizebox{1.0\textwidth}{!}{
\begin{threeparttable}
\input{table/reg_T02.tex}
\begin{tablenotes}
\footnotesize
\end{tablenotes}
\end{threeparttable}
}
\caption{Regression Results for Forecast Errors}
\label{tab:reg_fe}
\end{table}

Column (4) considers earnings volatility and stock price volatility as drivers of forecast error.
As predicted, both earnings and stock price volatility are positively correlated with forecast error, indicating analysts are less accurate about firms that have more volatile earnings growth and stock prices.
This result shows the robustness of the role of the number of analysts in driving the size of forecast error.
Relative to these results where we include all firms that have at least two analysts covering and reporting earnings forecasts, 
column (5) exclude firms that have only two analysts covering and reporting earnings forecasts, and the result is robust to doing so.
Column (6) excludes firms that have fewer than 6 analysts covering and reporting earnings forecasts.
The results remain robust with such sample selection, and it is confirmed that the number of analysts drives the size of forecast error.
In summary, the regression table examines factors associated with forecast error, finding that having more estimates reduces error, while greater variability in actual values and higher stock volatility are linked to larger forecast errors. The results are generally consistent across sub-sample specifications.





\subsection{Firms with higher earnings per share have more analysts}

We have shown that the number of analysts is a driver of forecast dispersion and forecast error measures.
Here, we explore what then drives the number of analysts for each firm.
Column (1) reports the result of regressing the number of analysts on the realized EPS, confirming the statistical significance of the relationship in Figure \ref{fig:sd_actual_numest}.
Column (2) adds year and firm dummies, and column (3) adds the standard other controls such as firm size and firm age.
As seen in column (3), larger and older firms have more analysts, and the positive relationship between the number of analysts and the realized EPS remains significant.
These three factors together with year and firm fixed effects yield a high $R^2$ of 0.844, which is mainly due to a high $R^2$ of 0.828 in column (2).

\begin{table}[htbp]
\centering
\resizebox{1.0\textwidth}{!}{
\begin{threeparttable}
\input{table/reg_T03.tex}
\begin{tablenotes}
\footnotesize
\end{tablenotes}
\end{threeparttable}
}
\caption{Regression Results for Number of Estimators}
\label{tab:reg_numest}
\end{table}


Column (4) considers earnings volatility and stock price volatility as drivers of the number of analysts.
The results thus far are about all firms that have at least two analysts covering and reporting earnings forecasts.
In column (5), we exclude firms that have only two analysts covering and reporting earnings forecasts, and the result is robust to doing so.
Column (5) shows that a one unit increase in realized EPS is associated with a 0.079 increase in the number of analysts.
Column (6) excludes firms that have fewer than 6 analysts covering and reporting earnings forecasts, showing that a one unit increase in realized EPS is associated with a 0.104 increase in the number of analysts, which is larger than the magnitude found in column (4) and column (5).
In summary, the regression table examines factors associated with the number of coverage analysts for each firm, and we find that there are more analysts for large and old firms, but it is also the case that realized EPS is a robust driver of the number of analysts across sub-sample specifications.



\bigskip
\section{Time-Series Index}


In preceeding sections, we have seen various characteristics of micro data based on IBES, NikkeiNeeds, and DBJ database.
The measure of forecast dispersion and forecast error appear to  reflect the level of imperfect information about a firm's future earnings.
In particular, the measure of forecastdispersion is a type of second momemnt proxy in terms of the level of imperfect  information, which can be 
thought of as the level of uncertainty about a firm's future earnings.
While it isa measure of such from the perspective of market analysts, not from the perspective of business managers, 
this section exploits the cross-sectional and time-series patterns of the data to construct time series indices of forecast dispersion and errors.

\begin{figure}[htbp]
    \centering
    \includegraphics[width=1.0\textwidth]{graph/meanfdiscv_m.png}
    \caption{Forecast Dispersion Index}
    \label{fig:fdis_full}
    \end{figure}

Figure \ref{fig:fdis_full} plots the forecast dispersion index, which is the cross-sectional mean of the forecast dispersion measure for each year.
We construct it by taking the annual mean of forecast dispersion for each firm and then the cross-sectional mean for each year.
As seen in Figure \ref{fig:fdis_full}, there are several spikes that are potentially corresponding to events that heightened uncertainty.
The most recent spike is the Covid-19 pandemic peak in March 2020.
The second highest peaks in the figure are from March 2003 when the Iraq War began, to May 2003 when the Japanese Government bailed out Resona Bank.
One may argue that these peaks are times when uncertainty facing firms in terms of their performance, such as earnings, is high.
Other spikes seen in the figure include when the Bank of Japan cut the  discount rate to 1.00 percent in 1994, 
during times when there was a banking crises in Japan, on top of theAsian financial crisis around 1998, 
the 9/11 terorist attack, the global financial crisis periods in 2007 through 2009, aftermath of Tohoku earthquake in 2013.

Here, we compare the forecast dispersion index with the Economic Policy Uncertainty (EPU) index for Japan, showing a positive correlation.
As seen in Figure \ref{fig:fdis_epu}, both the forecast dispersion index and EPU show some similarities in capturing times when uncertainty about firm performance appears to be high.


\begin{figure}[htbp]
    \centering
    \includegraphics[width=1.0\textwidth]{graph/mean_Fdis_CV_JPNEPUINDXM.png}
    \caption{Forecast Dispersion and Economic Policy Uncertainty in Japan}
    \label{fig:fdis_epu}
    \end{figure}
   

It is commonly understood that uncertainty is countercyclical in that it is higher when business conditions are worse, such as during recessions or when stock prices are falling.
To see whether the index created from IBES data has countercyclical features, Figure \ref{fig:fdis_nikkei} shows a relationship between the forecast dispersion index and the Nikkei 225 stock market index.
It can be seen that there are times when the forecast dispersion index rises and the Nikkei 225 stock market index falls, albeit with difficulty seeing it through with trends in the Nikkei 225 index.
While the jumps in the uncertainty index appear when firms face greater uncertainty about future earnings, it is also worthwhile to study how those jumps are clustered within some periods of time because
the frequency of those jumps seems to vary across time.

\begin{figure}[htbp]
    \centering
    \includegraphics[width=1.0\textwidth]{graph/mean_Fdis_CV_NIKKEI225.png}
    \caption{Forecast Dispersion and Nikkei 225 Index}
    \label{fig:fdis_nikkei}
    \end{figure}

    
Comparing the first half of the period in Figure 8 to the second half of the periods, it is clear that there are more jumps in the first half of the period than the second half. Hence, Figure 9 takes the moving average of the forecast dispersion index and forecast errors index as an effort to gauge the medium-run fluctuations of those uncertainty indexes. As one can see in Figure 9, it is evident that the overall level of uncertainty in the first half of the period is higher than that for the second half. As in the right panel of Figure 9, this pattern is more pronounced if you look at the forecast error index, where the level of uncertainty in the medium run is highest around the Asian financial crisis and the Japanese banking crisis before the uncertainty level fell sharply after the bailout of Resona Bank. Therefore, we argue that both the high-frequency time series of the uncertainty index with spikes and the moving average of such an uncertainty index are useful for researchers and policymakers to monitor and think about the level of uncertainty in real time.
The moving average of the forecast dispersion index and forecast error index (Figure \ref{fig:fdis_fe_ma}) provide a gauge of uncertainty over the medium term.

\begin{figure}[htbp]
\centering
\includegraphics[width=0.485\textwidth]{graph/mean_Fdis_CV_MA.png}
\includegraphics[width=0.485\textwidth]{graph/mean_FE_pct_MA.png}
\caption{Moving Average of Forecast Dispersion and Errors}
\label{fig:fdis_fe_ma}
\end{figure}


Finally, Table \ref{tab:cross_corr} reports the cross-correlations among the uncertainty indices and macroeconomic variables. The forecast dispersion and error indices are positively correlated with the EPU index and negatively correlated with the Nikkei 225 index.

\begin{table}[htbp]
\centering
\resizebox{1.0\textwidth}{!}{
\begin{threeparttable}
\input{table/cross_correlation.tex}
\begin{tablenotes}
\footnotesize
\item
\end{tablenotes}
\end{threeparttable}
}
\caption{Cross-Correlations}
\label{tab:cross_corr}
\end{table}

In summary, we construct the timeseries indices of forecast dispersion and forecast errors, and 
examine their basic  features in capturing times when uncertainty is high and compare them with the EPU index and the Nikkei 225 index.
We show that the forecast dispersion index and the forecast error index are positively correlated with the EPU index 
and negatively correlated with the Nikkei 225 index.
We argue that these indices are not only relevant for the level of uncertainty surrounding business performance in general but also 
useful for researchers and policymakers to monitor and think about the level of uncertainty in real time, as it is relatively easy 
to construct these indices as demonstrated in the paper.


\section{Conclusion}
In conclusion, our paper sheds new light on the properties anddeterminants of analyst forecast dispersion and errors, and their links to macroeconomic uncertainty. By constructing novel indices and exploring their cross-sectional and time-series patterns, we contribute to the growing literature on the role of expectations and uncertainty in the economy. Our findings suggest that analyst forecasts provide a valuable window into the evolution of firm-level uncertainty over time, with potentially important implications for researchers, policymakers, and market participants.
Our paper also points to several avenues for future research. One natural extension would be to examine the effects of forecast dispersion and errors on firm-level outcomes such as investment, employment, and asset prices. Another interesting question is whether the patterns we document for Japan also hold in other countries, particularly emerging markets where analyst coverage may be more limited. Finally, future work could explore the implications of our findings for the design of monetary and fiscal policies that aim to stabilize the economy in the face of uncertainty shocks.





\bibliographystyle{apalike}
\bibliography{ibes-japan}
\end{CJK}
\end{document}